%% notes-tex/publishing-system/sections/6-publishing.tex

\section{Publishing a build\secstatus{todo}}
\label{sec:publishing}

\notesec{Why Zotero}
The point of a review copy is annotation. Zotero's reader keeps highlights and
comments attached to a specific PDF attachment, so a comment always refers to the
exact rendering it was made on. A note saying a number is wrong is meaningless
unless the build it was made against can be identified, and a shared folder does
not give that.

\notesec{Where a build is filed}
The collection path is derived from the repository path, never configured:
\texttt{notes-tex/<slug>/} publishes into \texttt{torchcell / notes-tex / <slug>},
and \texttt{paper/nature-biotech/} into \texttt{torchcell / paper / nature-biotech}.
The parent collection is a flat index of every document of that kind and the leaf
holds the document's versions; the item is filed in both, because Zotero does not
show a sub-collection's items in its parent by default. The topic collections that
hold other people's papers are a separate tree, and keeping the two apart is what
stops a draft from being read as literature when a bibliography is rebuilt. Two
collections once shared the name \texttt{microbe-perturb-seq}, one for publishing
and one feeding \texttt{make bib}, and deriving the path from the directory is what
made that safe.

\notesec{How a version is named}
Each publish creates one child attachment named
\begin{center}
\texttt{<doc>\_<YYYY-MM-DD-HH-MM-SS>\_<sha256[:8]>.pdf}
\end{center}
The timestamp is what a
person sorts by; the hash identifies the bytes and is the same identifier the
repository's provenance rule uses everywhere else. A re-run of \texttt{make} with
nothing edited produces the same hash and is skipped as a duplicate, while any
real edit produces a new one. This rests on the reproducible build: without
\texttt{SOURCE\_DATE\_EPOCH} two builds of identical sources differed, and every
publish uploaded a fresh near-identical PDF. The git commit and branch the build
came from are written into the parent item, so a reviewer can go from a comment to
the exact source state.

\notesec{The commands}
For a note, \texttt{zotero\_publish.py <slug>} with \texttt{--dry-run} to preview,
\texttt{--list} to see the versions, \texttt{--clean} for the share view and
\texttt{--pdf <stem>} for any other PDF the document builds. For the manuscript,
\texttt{make publish} wraps the same script with \texttt{--pdf editing} and
\texttt{--tex sections/frontmatter.tex}, because the title is declared there and
the stock sample file in the same directory carries a placeholder title that a
search would also find. Credentials are the repository's \texttt{.env}.

\notesec{Landing}
A document is edited in a git worktree like any other change, and the commit
carries the built PDF and the converted figures beside the sources, so the PDF on
\texttt{main} is always the one its sources produce. The branch is opened as a
pull request and handed to the merge queue, which rebases, fast-forwards, closes
the pull request and removes the worktree; publishing to Zotero can happen before
or after, since the version is identified by content, not by commit.
